\section{Prapemrosesan dan Pembentukan Data Forecasting}
\label{sec:b3-prapemrosesan}

Prapemrosesan menyusun arsip observasi menjadi tensor masukan dan target yang memiliki urutan waktu serta urutan koridor yang konsisten. Urutan implementasinya adalah pemeriksaan timestamp dan penyelarasan grid, penyimpanan mask observasi, pengisian gap terisolasi, penentuan waktu valid dan segmen, pembagian kronologis, standardisasi berdasarkan data pelatihan, kemudian pembentukan window pada masing-masing pembagian data. Penanganan gap dan segmentasi dilakukan sebelum pembagian data; prosedur ini perlu dibedakan dari normalisasi yang statistiknya hanya dihitung pada data pelatihan.

\subsection{Pemeriksaan dan Penyelarasan Data Temporal}
\label{subsec:b3-penyelarasan}

Kolom waktu observasi dikonversi menjadi timestamp dengan zona waktu UTC. Catatan yang timestamp-nya tidak dapat dikonversi dibuang. Data setiap koridor kemudian diurutkan berdasarkan waktu, dan timestamp duplikat ditangani dengan mempertahankan catatan terakhir. Pemeriksaan ini menyelaraskan identitas waktu; keberadaan timestamp tidak dengan sendirinya menunjukkan bahwa nilai kecepatan tersedia.

Rentang bersama ditentukan menggunakan waktu mulai paling akhir dan waktu selesai paling awal di antara 64 koridor. Berdasarkan rentang tersebut dibentuk grid berinterval satu jam yang berisi 729 posisi waktu. Setiap deret koridor kemudian disejajarkan terhadap grid melalui reindexing. Proses ini tidak menghitung rata-rata kecepatan per jam dan tidak melakukan interpolasi; posisi yang tidak mempunyai catatan pada timestamp grid menjadi nilai hilang.

Sebelum imputasi, atribut is\_observed disimpan untuk setiap pasangan waktu--koridor. Nilainya benar apabila current\_speed tidak kosong pada posisi tersebut, dan salah apabila nilainya hilang. Mask ini menunjukkan ketersediaan kecepatan pada arsip, bukan hasil pengujian menyeluruh atas akurasi pengukuran TomTom. Sebanyak 24.967 dari 46.656 posisi waktu--koridor memiliki kecepatan terobservasi, sedangkan 21.689 posisi tidak tersedia sebelum imputasi. Angka tersebut menghitung pasangan waktu--koridor, bukan jumlah jam unik.

\begin{figure}[p]
\centering
\includegraphics[width=\linewidth,height=0.78\textheight,keepaspectratio]{assets/pics/bab3_03_missing_heatmap.png}
\caption[Ketersediaan observasi kecepatan sebelum imputasi]{Ketersediaan observasi kecepatan sebelum imputasi pada 64 koridor. Panel kanan memperbesar 48 jam pertama agar pola ketersediaan per jam terbaca. Sumber: mask is\_observed pada dataset penelitian, diverifikasi terhadap arsip mentah.}
\label{fig:b3-ketersediaan}
\end{figure}

Gambar~\ref{fig:b3-ketersediaan} menunjukkan posisi observasi asli dan nilai yang hilang. Pola kehilangan yang serupa antarkoridor perlu diperhatikan karena pemeriksaan gap pada pipeline dilakukan pada tingkat jaringan. Seleksi 64 koridor yang dijelaskan sebelumnya menggunakan jumlah timestamp pada seluruh arsip, termasuk periode setelah pelatihan. Dengan demikian, sampel koridor bersifat retrospektif dan hasilnya dibatasi pada jaringan terpilih. Hal ini tidak sama dengan memasukkan nilai target pengujian sebagai fitur, tetapi membatasi klaim generalisasi terhadap koridor baru atau pemilihan jaringan secara daring.

\subsection{Penanganan Missing Values dan Segmentasi Temporal}
\label{subsec:b3-imputasi}

Status kelengkapan pada suatu jam ditentukan dari mask kecepatan seluruh koridor. Suatu jam dinyatakan lengkap apabila semua koridor memiliki current\_speed terobservasi. Jam yang tidak lengkap menjadi kandidat gap satu jam terisolasi hanya jika jam tepat sebelumnya dan tepat sesudahnya sama-sama lengkap. Pada grid penelitian terdapat 320 jam yang memenuhi aturan kandidat tersebut. Jumlah ini bukan jumlah nilai yang diisi, sebab satu jam dapat mencakup lebih dari satu koridor.

Pada jam kandidat, nilai yang hilang pada current\_speed dan congestion\_ratio diisi menggunakan nilai atribut yang sama dari jam tepat sebelumnya pada koridor yang sama. Nilai yang sudah tersedia dipertahankan. Proses ini menggunakan pergeseran satu langkah dari deret sebelum pengisian, bukan pengisian berantai tanpa batas. Atribut was\_imputed ditandai benar apabila salah satu dari kedua atribut diisi; oleh karena itu, indikator tersebut tidak selalu berarti nilai kecepatanlah yang diimputasi. Mask is\_observed tetap mempertahankan status kecepatan sebelum pengisian.

Nilai pengganti berasal dari masa lalu sehingga tidak menggunakan kecepatan pada jam sesudahnya. Namun, keputusan apakah suatu gap boleh diisi menggunakan informasi ketersediaan observasi pada jam sesudahnya. Dengan demikian, aturan pemilihan gap bersifat retrospektif dan keseluruhan prosedur tidak dinyatakan sepenuhnya kausal untuk penggunaan daring. Masking target mencegah nilai kecepatan hasil imputasi diperlakukan sebagai target terobservasi, tetapi tidak menghilangkan keterbatasan aturan pemilihan gap ini.

Setelah pengisian, suatu timestamp dinyatakan valid untuk pembentukan matriks apabila current\_speed dan congestion\_ratio tidak kosong pada seluruh 64 koridor. Congestion\_ratio tidak menjadi fitur masukan model, tetapi ketidaktersediaannya dapat menggugurkan timestamp untuk seluruh jaringan. Kriteria tersebut memeriksa keberadaan nilai, bukan ambang confidence atau rentang kecepatan. Pemeriksaan ulang dataset menunjukkan bahwa pada artefak saat ini tidak ada timestamp tambahan yang gugur hanya akibat congestion\_ratio ketika seluruh kecepatan sudah tersedia.

Timestamp yang masih tidak valid tidak dimasukkan ke matriks pemodelan. Waktu valid yang berurutan dengan selisih tepat satu jam dikelompokkan dalam segmen yang sama; segmen baru dimulai ketika kesinambungan tersebut terputus. Dari 729 posisi waktu, sebanyak 710 waktu memenuhi kriteria validitas dan membentuk 10 segmen. Sebanyak 19 waktu yang tidak valid tetap dapat tercatat dalam berkas grid dengan identitas segmen kosong. Dengan pendekatan ini, jam yang hilang tidak dipadatkan seolah-olah dua observasi yang berjauhan merupakan dua jam berurutan.

\begin{figure}[htbp]
\centering
\includegraphics[width=\linewidth]{assets/pics/bab3_04_imputation_segments.png}
\caption[Contoh imputasi dan batas segmen temporal]{Contoh observasi asli, imputasi, dan batas segmen pada koridor Cawang, 21--22 Agustus 2026 (UTC). Periode dipilih di sekitar gap panjang pertama pada dataset, bukan berdasarkan hasil prediksi. Garis tidak disambungkan melintasi waktu yang tidak valid. Sumber: data olahan penelitian.}
\label{fig:b3-imputasi}
\end{figure}

Gambar~\ref{fig:b3-imputasi} memperlihatkan pengisian gap terisolasi dengan nilai sebelumnya serta pemutusan garis pada gap yang tidak dipertahankan. Penyajian kedua keadaan secara terpisah memperjelas perbedaan antara observasi yang tersedia, nilai pengganti, dan periode yang tidak digunakan dalam window.

\subsection{Pembagian Data Secara Kronologis}
\label{subsec:b3-split}

Pembagian data dilakukan berdasarkan urutan 729 posisi pada grid awal, bukan berdasarkan jumlah waktu valid atau jumlah window. Batas pelatihan ditetapkan pada $\lfloor0{,}70\times729\rfloor=510$ posisi, sedangkan batas akhir validasi pada $\lfloor0{,}85\times729\rfloor=619$ posisi. Pembagian tersebut menghasilkan 510 posisi pelatihan, 109 posisi validasi, dan 110 posisi pengujian. Tidak dilakukan pengacakan waktu untuk menentukan pembagian ini.

\begin{table}[htbp]
\centering
\caption{Pembagian kronologis pada grid dataset}
\label{tab:b3-split}
\renewcommand{\arraystretch}{1.15}
\begin{tabular}{>{\raggedright\arraybackslash}p{0.18\linewidth}>{\raggedright\arraybackslash}p{0.45\linewidth}rr}
\hline
\textbf{Bagian} & \textbf{Periode (UTC)} & \textbf{Grid} & \textbf{Valid}\\
\hline
Pelatihan & 12 Agustus 2026 04.00 -- 2 September 2026 09.00 & 510 & 491\\
Validasi & 2 September 2026 10.00 -- 6 September 2026 22.00 & 109 & 109\\
Pengujian & 6 September 2026 23.00 -- 11 September 2026 12.00 & 110 & 110\\
\hline
\end{tabular}
\end{table}

Pada Tabel~\ref{tab:b3-split}, kolom grid menyatakan jumlah posisi waktu sebelum penyaringan, sedangkan kolom valid menyatakan waktu yang memiliki identitas segmen. Pembentukan window dilakukan terpisah pada setiap bagian sehingga masukan dan target satu sampel tidak melintasi batas pelatihan, validasi, atau pengujian. Riwayat dari bagian sebelumnya tidak ditambahkan untuk membentuk window awal pada bagian berikutnya. Sebaliknya, penanganan gap sebelumnya dilakukan pada grid bersama; karena itu, pembentukan window terpisah tidak boleh ditafsirkan sebagai bukti bahwa seluruh keputusan prapemrosesan dilakukan secara independen per bagian.

\subsection{Normalisasi Data}
\label{subsec:b3-normalisasi}

Kecepatan dinormalisasi menggunakan standardisasi. Matriks kecepatan pelatihan pada timestamp valid diratakan menjadi satu kolom sebelum statistik dihitung. Dengan demikian, digunakan satu rata-rata dan satu simpangan baku bersama untuk seluruh koridor, bukan scaler terpisah per koridor. Sebanyak $491\times64=31.424$ nilai pelatihan digunakan, termasuk nilai yang telah diimputasi pada waktu valid. Nilai dari validasi dan pengujian tidak digunakan untuk menghitung statistik ini.

\begin{equation}
 z_{t,n}=\frac{v_{t,n}-\mu_{\mathrm{train}}}{\sigma_{\mathrm{train}}}.
 \label{eq:b3-standardisasi}
\end{equation}

Pada Persamaan~\ref{eq:b3-standardisasi}, $v_{t,n}$ adalah kecepatan koridor $n$ pada waktu $t$, sedangkan $z_{t,n}$ adalah nilai hasil standardisasi. Implementasi StandardScaler menggunakan simpangan baku populasi. Pemeriksaan ulang terhadap data pelatihan menghasilkan $\mu_{\mathrm{train}}=30{,}5066$ km/jam dan $\sigma_{\mathrm{train}}=11{,}3311$ km/jam, sesuai statistik yang tersimpan. Masukan dan target pada seluruh bagian ditransformasikan menggunakan statistik yang sama tanpa fit ulang pada validasi maupun pengujian. Pemisahan ini mencegah informasi distribusi kecepatan validasi dan pengujian masuk melalui estimasi scaler \cite{sklearn2026pitfalls}.

\begin{figure}[htbp]
\centering
\includegraphics[width=\linewidth]{assets/pics/bab3_05_split_normalization.png}
\caption[Alur pembagian data dan normalisasi]{Alur pembagian data dan normalisasi berdasarkan periode serta statistik dataset aktual. Lebar kotak tidak menyatakan proporsi durasi. Sumber: implementasi prapemrosesan dan scaler tersimpan.}
\label{fig:b3-split-normalisasi}
\end{figure}

Gambar~\ref{fig:b3-split-normalisasi} merangkum perbedaan antara penentuan waktu valid, pembagian grid, dan estimasi statistik pelatihan. Untuk interpretasi dalam km/jam, nilai yang telah distandardisasi dapat dikembalikan melalui $v=z\sigma_{\mathrm{train}}+\mu_{\mathrm{train}}$. Penggunaan scaler pelatihan saja tidak menghapus keterbatasan seleksi koridor dan pemilihan gap yang telah dijelaskan.

\subsection{Pembentukan Sliding Window dan Target Multi-Horizon}
\label{subsec:b3-window}

Matriks pada setiap bagian diurutkan berdasarkan waktu dan urutan node yang telah ditetapkan. Window digeser satu langkah untuk membentuk masukan sepanjang 12 langkah, yaitu dari $t-11$ sampai $t$. Target diambil pada $t+1$, $t+3$, dan $t+6$. Pada grid satu jam, jarak titik pertama dan terakhir masukan adalah 11 jam; istilah 12 langkah merujuk pada jumlah posisi observasinya. Ilustrasi hubungan masukan dan target mengikuti Gambar~\ref{fig:b2-window}.

Sampel hanya dipertahankan apabila seluruh indeks masukan dan ketiga target mempunyai identitas segmen yang sama. Karena segmen dibentuk dari waktu valid yang berurutan, aturan ini mencegah pengambilan target setelah gap panjang seolah-olah target tersebut berjarak satu, tiga, atau enam jam. Satu sampel mempunyai tensor masukan $X\in\mathbb{R}^{12\times64\times1}$ dan target $Y\in\mathbb{R}^{3\times64}$. Dimensi terakhir masukan hanya memuat kecepatan; mask observasi dan mask imputasi disimpan terpisah dan tidak ditambahkan sebagai kanal fitur.

Mask target diambil dari is\_observed pada waktu target, sehingga bentuknya juga $3\times64$. Nilai target yang berasal dari imputasi dapat tetap tersimpan dalam tensor, tetapi mask-nya bernilai salah dan tidak diperlakukan sebagai observasi aktual dalam evaluasi bermasker. Pembentukan dataset tidak membuang seluruh window hanya karena sebagian targetnya tidak terobservasi. Jumlah target valid dihitung sebagai pasangan sampel--koridor pada masing-masing horizon.

Hasil pembentukan dataset tersimpan memuat 415 sampel pelatihan, 92 sampel validasi, dan 93 sampel pengujian. Pada data pengujian terdapat 2.944 target valid untuk horizon satu jam, 2.944 untuk tiga jam, dan 3.008 untuk enam jam. Perbedaan himpunan target ini perlu diperhatikan saat membandingkan galat antarhorizon; perubahan galat tidak semata-mata menyatakan pengaruh jarak prediksi. Mask target, waktu target, serta mask masukan disimpan bersama tensor agar status observasi dapat ditelusuri saat evaluasi dan visualisasi.
